\section{La idea y el alcance del curso}

La teoría de esquemas es el lenguaje en el que la geometría algebraica moderna pone en correspondencia objetos geométricos y anillos conmutativos. La idea central, que guiará el curso, asocia a cada anillo $A$ su espectro primo $\Spec A$. Así, incluso $\Spec\mathbb Z$ se interpreta geométricamente y los problemas aritméticos pueden estudiarse con herramientas de geometría. Los puntos no serán únicamente soluciones sobre un cuerpo algebraicamente cerrado: serán todos los ideales primos, incluidos los puntos genéricos y los puntos que reflejan información aritmética.

La perspectiva aritmética estará presente desde el comienzo mediante ejemplos sobre $\mathbb Z$, cuerpos finitos, anillos locales y reducción módulo primos. El curso parte de los esquemas afines y su gavilla estructural; después construye esquemas generales por pegado y estudia sus morfismos y fibras. Más adelante aborda dimensión, regularidad, gavillas cuasicoherentes y coherentes, gavillas invertibles, divisores y la construcción $\Proj$. El programa concluye con curvas y ejemplos aritméticos. El programa académico distribuye estos temas en 24 sesiones y deja cohomología avanzada, planitud, geometría étale, descenso y stacks para un segundo curso.

El hilo conceptual inicial puede resumirse así: los ideales primos de $A$ son los puntos de $\Spec A$; la topología de Zariski describe sus especializaciones; las localizaciones $A_f$ y $A_{\mathfrak p}$ describen funciones en abiertos y gérmenes; la gavilla estructural reúne esa información local; y el pegado de abiertos afines produce esquemas que ya no tienen por qué ser afines.

\section{Lecturas y prerrequisitos}

Los textos principales indicados en el programa son \emph{The Rising Sea}, de Ravi Vakil, y \emph{Algebraic Geometry and Arithmetic Curves}, de Qing Liu. Como referencias complementarias aparecen Hartshorne, Görtz--Wedhorn, Mumford, Shafarevich, Atiyah--Macdonald, Eisenbud, el \emph{Stacks Project} y EGA. Para las primeras sesiones se espera familiaridad con ideales primos y maximales, localización, módulos, productos tensoriales, anillos noetherianos, dimensión de Krull, el lema de Nakayama y los teoremas de going-up y going-down; también se usa topología general básica. Las nociones de categorías, funtores, transformaciones naturales y propiedades universales se incorporan cuando hacen falta.

% \begin{nota}
% Este capítulo recoge la introducción temática del programa y de la presentación de la Sesión 0. Las indicaciones de inscripción, plataformas y evaluación pertenecen a la organización del curso y no forman parte del desarrollo matemático de los apuntes.
% \end{nota}
